\chapter{残った問い}\label{ch:open}

この章には、この研究で確かめられなかったことと、次に確かめる候補をまとめる。どの候補も、これまでと同じく、主な検定を 1 つに決めて事前登録してから 1 回だけ行う前提である。

\section{次の候補}

\begin{longtable}{P{0.6cm}P{4.2cm}P{5.4cm}P{4.2cm}}
\toprule
 & \textbf{候補} & \textbf{問い} & \textbf{必要なもの} \\
\midrule
\endhead
\bottomrule
\endfoot
1 & 文章と株価が食い違った発表 & 文章は良いと読めるのに day 0 の株価が下がった（またはその逆の）発表では、翌日以降に文章の向きへ戻る動きがあるか。市場が読み違えたなら戻るはず & 全部ある（実験 2 のパネル、実験 8 の翌日の反応）。CPU で数分 \\
2 & リターンではなくリスクを当てる & 発表文の読み取りは、過去の値動きでは分からない翌月のボラティリティを当てるか。当たれば、共分散の補正で TE の約束を守りやすくなる & ChronoGPT の埋め込み（ある）と、株価から作る目的変数 \\
3 & モデルどうしの意見の割れで $\Omega$ を決める & ChronoGPT・FinBERT・辞書・TF-IDF の点数の割れが大きいセクターほど、見通しを弱める。セクター単位の BLM が良くなるか & 実験 11 のセクター単位の BLM \\
4 & 中間の層の埋め込み & 実験 11 の探索で強かった 26 層目の埋め込みは、事前登録した比較でも従来の NLP と最後の層を超えるか & 埋め込み（ある） \\
5 & 10-K の文章と長い期間のリターン（実験 6） & 年次報告書のリスクの記述から作ったリスクへのさらされ方は、数か月〜1 年のリターンの横断面を説明するか & 10-K の取得と ChronoGPT の読み取り。事前登録の案はある（次の章） \\
6 & 同業他社への波及 & ある会社の発表の LLM の読み取りは、まだ発表していない同業他社の株価の反応も説明するか & 実験 2 のパネルとセクターの対応 \\
7 & 日中の反応 & 発表の直後の数分〜数時間で、LLM が市場より早く読めるか & 日中の株価（いまは持っていない） \\
\end{longtable}

候補 1〜3 は、2026 年 10 月 20 日の夜の議論で、この順に進めると決めた（\code{実験AとBの案\_議論用.md} の 6 章）。候補 1 が最も有望とした理由は、データがすべてそろっていて、問いが「市場はいつも正しく読むか」という、この研究の結論の核心に向いているからである。

\section{確かめられなかったこと}

\begin{description}[leftmargin=1.5em,style=nextline,itemsep=0.4em]
\item[gpt-4o-mini より強い LLM] gpt-4o-mini は小さく安いモデルである。より大きいモデルは反応をさらによく説明するかもしれない。ただし、知識の締め切りが新しいほど、評価できる期間（締め切りの後）は短くなる。
\item[アナリスト予想の売上とガイダンス] 実験 3 は EPS の予想だけを統制した。LLM の効果の一部は、売上の予想や、会社の見通しと市場の期待の差を文章から読んでいる可能性が残る。
\item[記憶の完全な排除] gpt-4o-mini の点数が先読みを含まないことは、ChronoGPT による傍証しかない。
\item[前向きの再現] 結果を見ていない将来の決算（たとえば 2026 年 10〜11 月の発表）で同じ検定をすることは、計画したが行わなかった（10 月 18 日の判断）。
\item[生存者バイアスの大きさ] 上場を廃止した銘柄の株価がない分の影響は、測っていない。
\item[追加のデータ] 内部者の取引、空売り、Wikipedia の閲覧数、マクロの系列は取得したが、実験には使っていない。
\end{description}

\section{LLM のシナリオと CVaR}

10 月 18 日に、「LLM に将来のシナリオを作らせ、モンテカルロで CVaR を計算してリスク管理に使う」研究を Gemini Deep Research で調べた（第 IX 部と付録）。照合できた文献から読めたのは、次の 3 点だった。
\begin{enumerate}
\item LLM にマクロのショックの筋書きを JSON で出させ、ありえない値を規則で落としてから、ファクター・モデルで損失に変える研究はある。ただし検証は「数値が安定する」までで、シナリオが当たるかは確かめられていない。
\item LLM の数値の予想は楽観に偏り、自信過剰である。下落の大きさや確率は LLM に任せないほうがよい（この研究の「回答の 76\% が正」と同じ向き）。
\item シナリオを入れる器としては、BLM より、過去のシナリオの重みだけを動かす Entropy Pooling（Meucci）が向いている。
\end{enumerate}
「LLM に筋書き、データに大きさと確率」という分担は、この研究の考え方と同じである。この研究の結論（LLM の読み取りは期待リターンでは価格に入りきる）からも、候補 2（リスクを当てる）と同じく、LLM をリスクの側に使う方向は検討に値する。ただし 2023-11 以降には大きな下落がほとんどなく、CVaR の検証の検出力は低い。
